\paragraph{Information convention and comparison.}
The Fourier representation used below is the established Bayesian
representation discussed in Section~\ref{sec:fourier-filtering-comparison}.
The resource restriction is instead a persistent index in an alphabet
of size at most $M$. The prior, contrast, horizon, and interface may
enter a read-only program; a transition receives only its preceding
index and the current command and report. In particular, an exact
Fourier vector is not charged as one label. The causal assertion of
Theorem~\ref{thm:circle-resolution} uses this convention, at fixed
$N$ and known contrast; it is not an effective numerical compiler.

For orientation, consider one past trial and one future trial. If the
past factor is $a\{1+\tau(zw+\bar z w^{-1})\}$, Haar normalization
gives $c_1=\tau z$. With the query parameter $\rho$ and phases
$\varphi_l$ defined above, direct multiplication gives
\[
 q_l(z)=\tfrac12+\rho\tau^2\Re(ze^{i\varphi_l}),\qquad
 \frac13\sum_{l=0}^{2}|q_l(z)-q_l(z')|^2
       =\tfrac12\rho^2\tau^4|z-z'|^2.
\]
The two contrast factors are already visible in this elementary case.
The higher-dimensional result requires more: the normalized flags must
be excited by a set of acquisition histories with uniformly positive
subprobability, the global image must admit a matching cover, and the
successive covers must be compatible with the weighted Bayesian update.
The next two lemmas and the proof of the theorem establish these
properties with the actual evidence retained. This computation is an
illustration of the theorem, not a separate improvement of its law.
